\section{Actual removal of isolated CI cancellation by tuple normalization}
\label{sec:r9-isolated-ci-cancellation}

The frozen CI0 theorem requires the support of the difference
to be the union of the two zero-scheme supports.  We now remove
that restriction for a checkable class of separated completely
cancelled components.  We change generating tuples, not zero
schemes, the original cycle, its exact complement, the line
bundle, or the metric.  The proof gives an actual form directly
and requires no compact-curvature solver.

\subsection{Canonical conormal comparison}

We first give the local algebra used to formulate the hypothesis.
If $f_1,\ldots,f_r$ is a regular sequence in a local ring $R$,
every relation $\sum_i a_i f_i=0$ has all coefficients in
$I=(f_1,\ldots,f_r)$.  Here is the required elementary argument.
For $r=1$ the non-zero-divisor condition gives $a_1=0$.
For the induction step, reduce modulo $(f_1,\ldots,f_{r-1})$.
The last regularity condition gives
$a_r=\sum_{i<r}b_i f_i$.  Substitution leaves the relation
$\sum_{i<r}(a_i+b_i f_r)f_i=0$ among the first $r-1$
generators, to which induction applies.  This proves the claim.

Consequently, if two $r$-tuples $f,g$ generate the same ideal
$I$ near their common zero locus and $f$ is regular, a matrix
$A$ with $g=Af$ is unique modulo $I$.  A matrix $B$ with
$f=Bg$ satisfies $(BA-I)f=0$, so $BA\equiv I\pmod I$.
The resulting square matrix $\overline A$ is invertible over
$R/I$, and $A$ is invertible after shrinking about the zero
locus.  These conclusions do not require the common scheme
to be reduced.

Let now $X/K$ be smooth projective of dimension $n$, let
$1\le r\le n$, and let $s_0,s_1$ be regular $r$-tuples of
global sections of the same line bundle $L$.  Write their pure
codimension-$r$ zero schemes as $Y_0,Y_1$.  Let $C\subset X$
be a reduced closed subscheme on a neighborhood of which the
two zero schemes agree.  In a common local trivialization
write $s_1=A_\nu s_0$.  The matrices $A_\nu|_C$ agree on
overlaps by the relation argument, and changing the frame of
$L$ does not change the comparison matrix.  They therefore
define a canonical global matrix
\[
         A_C\in\mathrm{GL}_r\bigl(H^0(C,\mathcal O_C)\bigr).
 \tag{ICC1}
\]
Equivalently it is the comparison between the two displayed
framings of the common conormal sheaf, restricted to $C$.
If $H^0(C,\mathcal O_C)=K$, then $A_C$ is a literal constant
matrix in $\mathrm{GL}_r(K)$.  Equality of cycle multiplicities
does not imply the equality of ideals used here.

\subsection{A finite, checkable cancellation theorem}

Put
\[
 Z=[Y_1]-[Y_0],\quad q=n-r,\quad
 U=X^{\rm an}\setminus|Z|,\quad
 V=X^{\rm an}\setminus(|Y_0|\cup|Y_1|).
 \tag{ICC2}
\]
Assume the following algebraic and finite-matrix conditions.
There are pairwise disjoint reduced proper closed subschemes
$C_1,\ldots,C_t$, each with $H^0(C_j,\mathcal O_{C_j})=K$,
such that
\[
 |Y_0|\cup|Y_1|=|Z|\ \amalg\ \coprod_{j=1}^t|C_j|,
 \tag{ICC3}
\]
and the ideals of $Y_0,Y_1$ agree on a Zariski neighborhood
of every $C_j$.  Let $A_j\in\mathrm{GL}_r(K)$ be (ICC1),
choose one index, denoted $1$, and assume
\[
                  A_1^{-1}A_j\in\mathrm{GL}_r(\mathcal O_K)
                              \quad(1\le j\le t).
 \tag{ICC4}
\]
This is a finite valuation test: every entry of each displayed
matrix and its inverse must have norm at most one.  When
$t=1$ there is no matrix condition to check.

\begin{theorem}[Isolated cancellation with an integral conormal certificate]
\label{thm:r9-isolated-ci-cancellation}
Under (ICC3)--(ICC4), replace the second tuple globally by
\[
                       \widetilde s_1=A_1^{-1}s_1.
 \tag{ICC5}
\]
For any one piecewise smooth metric on $L$, let
$G_0,G_{\widetilde1}$ be the actual frozen CI2 expressions.
Their difference on $V$ extends by zero across all $C_j$ to
an actual form
\[
                    g_Z\in B^{r-1,r-1}(U)
 \tag{ICC6}
\]
over the original $K$.  For every actual compact test,
\[
       \int_U g_Z\wedge D\eta=-\int_Z\eta,
                    \qquad\eta\in B_c^{q,q}(X^{\rm an}).
 \tag{ICC7}
\]
In particular a single completely cancelled $K$-rational
point, with identical local zero-scheme ideals, admits this
construction in every dimension $n=r$.  More generally it
works in every $n\ge r$ for a separated common component
with the stated constant global functions, even when its
common scheme structure is nonreduced.
\end{theorem}

\begin{proof}
An invertible constant matrix changes no ideal generated by
the components of a global tuple.  Thus (ICC5) leaves the
scheme $Y_1$ and all its generic lengths unchanged.  Its
components remain global sections of exactly the same $L$.

Near a point of $C_j$, choose a local matrix $A_\nu$ with
$s_1=A_\nu s_0$.  It restricts to $A_j$ on $C_j$.  Put
$Q_j=A_1^{-1}A_j\in\mathrm{GL}_r(\mathcal O_K)$ and
\[
            E_\nu=Q_j^{-1}A_1^{-1}A_\nu-I.
 \tag{ICC8}
\]
All entries of $E_\nu$ vanish on $C_j$.  Shrink to the
analytic open set on which every entry has norm less than
one.  For every Berkovich point $x$ in it and every vector
$v\in\mathscr H(x)^r$, the ultrametric isometry calculation
and the integrality of $Q_j,Q_j^{-1}$ give
\[
                 \|A_1^{-1}A_\nu(x)v\|_\infty
                                      =\|v\|_\infty.
 \tag{ICC9}
\]
Apply this to the value vector of $s_0$ in the common local
frame.  The common metric contributes the same scalar to
each logarithmic minimum, so the actual potentials satisfy
$\varphi_{\widetilde1}=\varphi_0$ on the punctured open set.
Consequently $\omega-D\varphi_{\widetilde1}
=\omega-D\varphi_0$, and the entire CI2 expressions agree
literally as actual forms, including every mixed metric term.

These analytic open sets cover $C_j^{\rm an}$.  They can be
intersected with the complement of $|Z|$ by (ICC3); denote
their union by $N_j$.  On $N_j\cap V$ the two forms coincide.
The zero form on $N_j$ therefore glues with
$G_{\widetilde1}-G_0$ on $V$.  Since
$U=V\cup\bigcup_jN_j$, the actual sheaf property gives
(ICC6) on the original exact complement.  This is actual
zero extension on an open cover, not extension of a current.

For $\eta\in B_c^{q,q}(X^{\rm an})$, the differentiated test
$D\eta$ has degree $(q+1,q+1)$.  The frozen actual support
lemma therefore makes its support compactly disjoint from
the entire union $|Y_0|\cup|Y_1|$, not just from $|Z|$.
Hence
\[
 \int_Ug_Z\wedge D\eta
     =\int_V(G_{\widetilde1}-G_0)\wedge D\eta.
 \tag{ICC10}
\]
Apply the frozen CI8 all-actual identity to each tuple.
Their zero schemes are the original $Y_0,Y_1$ and their
curvature terms are the identical actual form $\omega^r$.
Subtracting gives exactly (ICC7), with the original cycle
and its original multiplicities.
\end{proof}

For each fixed compact target presentation in degree
$(q+1,q+1)$ the same support lemma supplies one common
neighborhood of the entire union of the zero supports on
which all of its elements vanish.  On the compact complement
the fixed form (ICC6) has a finite actual presentation.
Common refinement, wedge, and integration therefore give
the same finite mixed coefficient seminorm bounds as in the
frozen CI proof.  Thus $I(g_Z)$ is continuous for the full
coefficient-final topology.  This continuity is a consequence
of the actual form just constructed, not its substitute.

\subsection{Exact scope of the constant normalization condition}

Condition (ICC4) is also necessary for a single left constant
normalization to make all the canonical conormal matrices
universal sup-norm isometries.  Indeed, if
$MA_j\in\mathrm{GL}_r(\mathcal O_K)$ for every $j$, then
\[
          A_1^{-1}A_j=(MA_1)^{-1}(MA_j)
                              \in\mathrm{GL}_r(\mathcal O_K).
 \tag{ICC11}
\]
The converse is (ICC5).  This is necessity for that precise
normalization procedure, not necessity for actual extension
of the CI expression for a particular tuple and not necessity
for strong G.  Changing both tuples, changing other allowed
data, or constructing another actual solution may succeed.

For example $A_1=I$ and
$A_2=\operatorname{diag}(\pi,\pi^{-1},1,\ldots,1)$, where
$0<|\pi|<1$, have determinant of norm one but fail (ICC4).
Thus determinant valuation alone does not certify simultaneous
normalization.  This is a finite-matrix route failure, not
an actual-cycle counterexample.

\subsection{Calibration on the frozen CNC family}

Take exactly the R8 CNC data, for every $n\ge r\ge2$:
\[
 \begin{gathered}
 X=\mathbf P^r\times(\mathbf P^1)^{n-r},\quad
 L=\operatorname{pr}_1^*\mathcal O(2),\quad
 s_0=(X_0X_1,X_2^2,\ldots,X_r^2),\\
 s_1=(aX_1(X_0-X_1),X_2^2,\ldots,X_r^2).
 \end{gathered}
 \tag{ICC12}
\]
The common cancelled support is
$C=[1:0:\cdots:0]\times(\mathbf P^1)^{n-r}$.
On $X_0\ne0$, put $z_i=X_i/X_0$.  Near $C$ the two tuples
are $(z_1,z_2^2,\ldots,z_r^2)$ and
$(a(1-z_1)z_1,z_2^2,\ldots,z_r^2)$ in their common frame.
Their ideals agree where $1-z_1$ is invertible, including
all the nonreduced transverse square structure, and
\[
                      A_C=\operatorname{diag}(a,1,\ldots,1).
 \tag{ICC13}
\]
Global functions on the product of projective lines are $K$:
restrict a regular function to each projective-line fiber,
where it is constant, and iterate through the factors.
Thus the theorem applies with $t=1$.  Its normalized tuple is
\[
                  \widetilde s_1=(X_1(X_0-X_1),X_2^2,\ldots,X_r^2).
 \tag{ICC14}
\]
On the original-field analytic open $|z_1|<1$ one has
$|1-z_1|=1$, so the two maxima, their metric potentials, and
their entire actual CI forms agree exactly.  The difference
therefore extends by zero through $C$ and gives (ICC7) for
the original nonzero cycle
$2^{r-1}([Z_1]-[Z_0])$ and original exact complement.
This uses the same $\mathcal O(2)$ and the same metric as the
specified expression that failed to extend in R8.

There is no contradiction with CNC14: the unnormalized
$G_{s_1}-G_{s_0}$ remains nonextendable.  The correction
\[
                 k=G_{s_1}-G_{\widetilde s_1}
                         \quad\text{on }X^{\rm an}\setminus|Y_1|
 \tag{ICC15}
\]
is a literal actual form.  Since both tuples have the identical
zero scheme and metric, subtraction of their CI8 identities
shows that it annihilates every $D\eta$ of the required degree.
Subtracting this explicitly constructed correction from the
old specified difference gives the new form (ICC6).
R8 already proved existence for this particular cycle via
another line bundle; the new assertion is the general
normalization theorem and the same-line-bundle cancellation
mechanism, not a new existence claim for that old example.

\paragraph{Unresolved interfaces.}
No common ideal is inferred merely from equal cycle lengths.
No constancy of the conormal comparison is inferred when
$H^0(C,\mathcal O_C)\ne K$.
No extension-degree correction or arbitrary actual descent
is used.  For components meeting the uncancelled support,
or incompatible conormal matrices, the theorem does not
provide a construction.  It solves neither the full
lower-middle OPEN-KERNEL problem nor general strong G,
and it supplies none of the further independence, symmetry,
products, or global gluing required for a complete BB height.
